% ============================================================================
% Capítulo 6 — Arquitetura de software
% ============================================================================
\section{ARQUITETURA DE SOFTWARE}

\subsection{Estilo arquitetural}

A solução adota o estilo de \textbf{monólito modular em camadas}, com separação
rígida de dependências: o domínio não conhece infraestrutura, a aplicação
depende apenas do domínio e a infraestrutura implementa as abstrações definidas
nas camadas internas. Essa organização segue os princípios de inversão de
dependência e permite substituir detalhes tecnológicos -- banco, armazenamento de
arquivos, provedor de tokens -- sem alterar as regras de negócio (MARTIN, 2017;
EVANS, 2003).

A escolha por um monólito modular, em vez de microsserviços, é deliberada. O
volume de operações da cooperativa não justifica a complexidade operacional de
serviços distribuídos, e a equipe reduzida teria dificuldade de sustentar
múltiplas implantações. O monólito modular preserva a possibilidade de extração
futura de módulos, caso o crescimento o exija, sem pagar antecipadamente o custo
de coordenação distribuída.

\begin{figure}[H]
  \centering
  \caption{Arquitetura em camadas da solução (a seta à direita representa a inversão de dependência)}
  \label{fig:camadas}
  \begin{tikzpicture}[
    camada/.style={draw, rounded corners=3pt, minimum width=130mm, minimum height=13mm,
      align=center, font=\fontsize{9.5}{11}\selectfont},
    seta/.style={-latex, thick, gray!70}
  ]
    \node[camada, fill=blue!8] (apres) {\textbf{Apresentação} --- MVC, Web API, Swagger, views Razor};
    \node[camada, fill=green!8, below=3mm of apres] (aplic) {\textbf{Aplicação} --- casos de uso, validações, DTOs, autorização};
    \node[camada, fill=orange!10, below=3mm of aplic] (dom) {\textbf{Domínio} --- entidades, enums, invariantes, contratos de repositório};
    \node[camada, fill=gray!10, below=3mm of dom] (infra) {\textbf{Infraestrutura} --- EF Core, SQL Server, JWT, arquivos, seed};
    \draw[seta] (apres) -- (aplic);
    \draw[seta] (aplic) -- (dom);
    \draw[seta] (infra.east) -- ++(1.4,0) |- (dom.east);
  \end{tikzpicture}
  \fonte{elaborado pelo autor (2026)}
\end{figure}

\subsection{Camadas e responsabilidades}

A Tabela~\ref{tab:camadas} detalha o papel de cada projeto da solução e os
artefatos que concentra.

\begin{table}[H]
  \centering
  \caption{Camadas, projetos e responsabilidades}
  \label{tab:camadas}
  \fontsize{9.5}{11.5}\selectfont
  \begin{tabular}{L{2.4cm}L{4.5cm}L{7.9cm}}
    \toprule
    \textbf{Camada} & \textbf{Projeto} & \textbf{Responsabilidade} \\
    \midrule
    Domínio & \texttt{AgroConecta.Domain} & Entidades, enums, invariantes,
      contratos de repositório e unidade de trabalho \\
    Aplicação & \texttt{AgroConecta.Application} & Casos de uso, comandos,
      consultas, validações e regras de orquestração \\
    Infraestrutura & \texttt{AgroConecta.Infrastructure} & Contexto EF Core,
      mapeamentos, repositórios, migrations, seed e serviços técnicos \\
    Apresentação & \texttt{AgroConecta.Web} & Controladores MVC, API REST,
      views, filtros, middlewares e configuração \\
    Testes & \texttt{AgroConecta.Tests} & Testes unitários, de integração e de
      API \\
    \bottomrule
  \end{tabular}

  \vspace{2pt}
  {\fontsize{10}{12}\selectfont\itshape Fonte: elaborado pelo autor (2026)}
\end{table}

\subsection{Orientação a objetos e princípios SOLID}

O modelo de domínio utiliza classes com estado protegido e comportamento
explícito. Todas as entidades herdam de \texttt{EntidadeBase}, que concentra
identificador, datas de criação e atualização e a comparação por identidade --
evitando duplicação e garantindo consistência transversal. As transições de
estado são implementadas como métodos, e não como atribuições diretas: uma
entrega só pode ser marcada como concluída se estiver em rota, e um pedido só
pode ser cancelado enquanto não houver fatura emitida. Essa abordagem mantém as
invariantes próximas dos dados que protegem.

Os princípios SOLID estão presentes de forma concreta: a
\textbf{responsabilidade única} separa entidade, caso de uso e repositório; o
\textbf{aberto/fechado} permitiu incluir rotas e faturamento sem alterar
entidades existentes; a \textbf{segregação de interfaces} mantém um contrato por
agregado; e a \textbf{inversão de dependência} sustenta a arquitetura, com a
aplicação declarando o que precisa e a infraestrutura fornecendo a
implementação.

\subsection{Padrões aplicados}

A Tabela~\ref{tab:padroes} resume os padrões utilizados e sua justificativa no
contexto do projeto.

\begin{table}[H]
  \centering
  \caption{Padrões de projeto aplicados}
  \label{tab:padroes}
  \fontsize{9.5}{11.5}\selectfont
  \begin{tabular}{L{3.4cm}L{5.8cm}L{5.8cm}}
    \toprule
    \textbf{Padrão} & \textbf{Aplicação} & \textbf{Benefício} \\
    \midrule
    Repository & Interfaces no domínio, implementação com EF Core &
      Isola a persistência e facilita testes com dublês \\
    Unit of Work & \texttt{IUnidadeDeTrabalho} com transação única por caso de
      uso & Consistência entre agregados \\
    Caso de uso & Uma classe por operação (\texttt{CriarPedido},
      \texttt{GerarFatura}) & Regras explícitas e testáveis \\
    DTO & Contratos de entrada e saída separados das entidades & Evita
      exposição do modelo interno \\
    Result & Retorno explícito de sucesso ou falha de negócio & Erros
      previsíveis sem exceções de controle de fluxo \\
    Problem Details & Respostas de erro padronizadas (RFC 7807) & Integração
      simples para o cliente mobile \\
    \bottomrule
  \end{tabular}

  \vspace{2pt}
  {\fontsize{10}{12}\selectfont\itshape Fonte: elaborado pelo autor (2026)}
\end{table}

\subsection{Modularização e dependências}

A solução é organizada em uma solução única com seis projetos, cujas referências
apontam sempre para dentro: \texttt{Web} depende de \texttt{Application} e
\texttt{Infrastructure}; \texttt{Infrastructure} depende de
\texttt{Application} e \texttt{Domain}; \texttt{Application} depende apenas de
\texttt{Domain}; e \texttt{Domain} não depende de nenhum outro projeto. Os
testes referenciam todas as camadas, mas não são referenciados por nenhuma.

Essa disciplina produz dois efeitos práticos. Primeiro, impede que detalhes de
banco contaminem as regras de negócio. Segundo, torna o domínio e a aplicação
plenamente testáveis sem servidor ou banco, o que explica a cobertura de 92,5\%
e 93,8\% nessas camadas, respectivamente, com 159 testes automatizados. A
verificação de arquitetura é feita pelo compilador: qualquer referência
indevida quebra o build, e não apenas uma convenção de equipe.
